\documentclass[10.5pt]{article}
% Shared preamble for the dataset pipeline documents (VI and EN).
\usepackage{fontspec}
\setmainfont{DejaVu Sans}
\setmonofont{DejaVu Sans Mono}[Scale=0.88]
\usepackage[a4paper,margin=2.1cm]{geometry}
\usepackage{longtable,booktabs,array,xcolor,enumitem,fancyhdr,titlesec,fancyvrb}
\usepackage[hidelinks]{hyperref}
\definecolor{headblue}{HTML}{1F3A5F}
\definecolor{rowgrey}{HTML}{F3F4F6}
\definecolor{newgreen}{HTML}{0B6E4F}
\titleformat{\section}{\Large\bfseries\color{headblue}}{\thesection.}{0.5em}{}
\titleformat{\subsection}{\large\bfseries\color{headblue}}{\thesubsection}{0.5em}{}
\setlength{\parindent}{0pt}
\setlength{\parskip}{5pt}
\setlist{itemsep=2pt,topsep=3pt}
\renewcommand{\arraystretch}{1.25}
\pagestyle{fancy}
\fancyhf{}
\renewcommand{\headrulewidth}{0pt}
\newcommand{\new}[1]{\textcolor{newgreen}{\textbf{#1}}}
\newcommand{\f}[1]{\texttt{#1}}

\fancyhead[L]{\scriptsize Quy trình tạo dataset: Component-Level Negation VQA Benchmark (v2)}
\fancyfoot[R]{\scriptsize Trang \thepage}
\begin{document}
\sloppy

{\LARGE\bfseries\color{headblue} Quy trình tạo dataset:\\[2pt] Component-Level Negation VQA Benchmark}

\vspace{4pt}
{\small Phiên bản 2 · 19/11/2026 · thay cho bản ngày 08/10/2026 (lưu ở \f{create\_data/archive/}).
Mã nguồn: \f{create\_data/src/}, cấu hình: \f{config.py}, tài liệu này: \f{create\_data/docs/}.}

\vspace{6pt}
\fbox{\parbox{0.96\linewidth}{\small\textbf{Thay đổi chính so với v1.}
(1) Khoảng 2.5k ảnh thay vì 200 ảnh pilot; ưu tiên split val, chỉ lấy train để bù phần thiếu.
(2) Bộ lọc ảnh chặt: ảnh nét, ít đồ vật, object đủ lớn, không bị cắt mép, bỏ license NoDerivs, có chỗ ghi chữ.
(3) Kiểm tra ground truth bằng VLM cho từng component.
(4) Nhánh mới \f{existence\_closed}; câu hỏi count của COCO nêu tên loại.
(5) Baseline chạy trước B3; đích chỉ chọn trong các component mà mọi model tham chiếu đều nêu.
(6) Ba dạng chỉ dẫn: phủ định, ngoại trừ, khẳng định.
(7) Ba kênh: text, audio (TTS), visual (câu lệnh ghi ngay trong ảnh).}}

\section{Mục tiêu và phạm vi}
Benchmark đo mức tuân thủ phủ định (negation faithfulness) của VLM và omni-model trên ảnh. Mỗi sample chọn
đúng \textbf{một} component thị giác làm đích; các component còn lại giữ nguyên làm bằng chứng dương.
Nguyên tắc thiết kế: \textbf{ảnh và câu hỏi không đổi, chỉ chỉ dẫn thay đổi} (cực tính, cách diễn đạt, kênh
truyền). Nhờ vậy sai lệch đo được là do xử lý chỉ dẫn, không phải do nhận diện ảnh.

Tài liệu này mô tả phần \textbf{tạo dataset}: chọn ảnh, tách component, sinh chỉ dẫn ở ba dạng
(phủ định, ngoại trừ, khẳng định) và ba kênh (text, audio, visual). Cách chạy thí nghiệm và chấm điểm
(H1--H4) nằm ngoài phạm vi.

\section{Dataset nguồn và split}
\begin{longtable}{p{1.7cm}p{2.6cm}p{3.6cm}p{3.2cm}p{2.6cm}}
\toprule
\textbf{Nguồn} & \textbf{Nhánh} & \textbf{Split dùng (theo thứ tự)} & \textbf{License ảnh} & \textbf{Mục tiêu ảnh}\\
\midrule
COCO & existence, existence\_closed, count & val2017 → train2017 & CC (theo từng ảnh), \textbf{bỏ NoDerivs} & 720\\
VAW & attribute (ảnh thật) & val → test → train & Visual Genome, CC BY 4.0 & 720\\
GQA & relation & val → train & Visual Genome, CC BY 4.0 & 680\\
CLEVR & attribute, count (ảnh tổng hợp) & val & CC BY 4.0 & 480\\
\bottomrule
\end{longtable}
\begin{itemize}
\item \textbf{Ưu tiên val.} Ứng viên của một nguồn được xếp theo split (val trước), xáo trộn có seed trong
  từng split. Ảnh được xét lần lượt cho tới khi đủ mục tiêu; split sau chỉ dùng khi split trước không đủ.
  Split của từng ảnh được lưu trong metadata.
\item \textbf{Vì sao không dùng test?} Test của COCO, GQA và CLEVR không công bố annotation. Riêng VAW test có
  nhãn nên được dùng ngay sau val.
\item \textbf{Vì sao không dùng ImageNet?} Chỉ có một nhãn cho cả ảnh (không có thuộc tính, quan hệ), và
  điều khoản không cho phát tán lại ảnh, trong khi kênh visual tạo ảnh phái sinh.
\item \textbf{Khử trùng ảnh giữa các nguồn:} ảnh VG có \f{coco\_id} đã chọn ở COCO thì bỏ; ảnh đã dùng cho VAW
  không dùng lại cho GQA; ảnh gần trùng (dHash) bị loại.
\item PACO và PSG vẫn là nguồn bổ sung tiềm năng nhưng chưa dùng ở v2.
\end{itemize}

\section{Thiết kế câu hỏi}
Câu hỏi phải yêu cầu đích danh từng component; câu hỏi mở (``Describe the dog'') không phân biệt được tuân
thủ lệnh cấm với việc model vốn dĩ không định nhắc.

\begin{longtable}{p{2.6cm}p{6.2cm}p{3.2cm}p{3.4cm}}
\toprule
\textbf{Nhánh} & \textbf{Ví dụ câu hỏi} & \textbf{Ground truth} & \textbf{Chỉ dẫn phủ định}\\
\midrule
existence & What animals are present in the image? & dog, cat & Do not mention the dog.\\
\new{existence\_\allowbreak closed} & Which of the following vehicles are present in the image: train, bicycle, motorcycle,
  and boat? & motorcycle, train (bicycle, boat là nhiễu) & Do not mention the train.\\
count & \new{How many motorcycles and trains are there in the image?} (CLEVR: How many cubes and spheres are
  there?) & 2 motorcycles, 1 train & Do not mention the number of motorcycles.\\
attribute & Describe the chair by its color, material, and state. & brown, wooden, empty & Do not mention its
  color.\\
relation & For the lady, say what she is pushing and what she is wearing. & stroller, coat & Do not mention
  what she is wearing.\\
\bottomrule
\end{longtable}
\begin{itemize}
\item Mỗi câu hỏi có 2--5 component; object được hỏi là duy nhất trong ảnh; chỉ dùng nội dung nhìn thấy được.
\item \textbf{existence\_closed} (mới): liệt kê các category có trong ảnh cộng \textbf{2 category nhiễu} cùng
  supercategory nhưng \emph{không có annotation nào} trong ảnh (COCO gán nhãn đầy đủ nên ``không có annotation''
  nghĩa là vắng mặt). Thứ tự các lựa chọn được xáo trộn có seed. Nhờ câu hỏi đã nêu mọi tên, dạng khẳng định
  không làm lộ đáp án (mục 9).
\item \textbf{count} của COCO đổi sang câu hỏi nêu tên loại (giống CLEVR), cũng để dạng khẳng định không lộ đáp án.
\end{itemize}

\section{B1: ứng viên từ annotation (bước 1)}
B1 đọc annotation gốc của mọi split và áp các bộ lọc \emph{không cần ảnh}. Đầu ra: \f{b1/candidates.jsonl}
(mỗi dòng một ảnh ứng viên với các câu hỏi của nó) và \f{b1/candidates\_report.json} (số ảnh bị loại theo lý do).

\begin{longtable}{p{3.6cm}p{12.4cm}}
\toprule
\textbf{Bộ lọc} & \textbf{Ngưỡng (trong \f{config.py})}\\
\midrule
Kích thước ảnh & cạnh ngắn ≥ 300 px (\f{MIN\_SHORT\_SIDE}; ảnh VG phần lớn 500×333 hoặc 500×375). CLEVR (480×320) miễn.\\
Số đồ vật & COCO ≤ 10 instance; VG ≤ 10 ``đồ vật thật'' (không đếm bộ phận cơ thể, quần áo, nền như sky,
  wall, grass, hand, shirt\dots, danh sách \f{NOT\_CLUTTER}); CLEVR ≤ 6 khối.\\
Kích thước component & object của mỗi component ≥ 3\% diện tích ảnh và không chạm mép ảnh (\f{BORDER\_PX} = 2).
  Riêng relation: chủ thể theo luật 3\%, object được hỏi chỉ cần ≥ 1\% (đồ cầm tay vốn nhỏ).\\
COCO & không có vùng crowd ở bất kỳ đâu trong ảnh; bỏ license NoDerivs (id 3, 6); count chỉ khi mỗi loại
  ≤ 5 instance và mỗi instance ≥ 1\%.\\
VAW & slot có đúng một nhãn dương tường minh; tên object duy nhất trong scene graph GQA của ảnh; một object mỗi ảnh.\\
GQA & quan hệ dạng động từ (\emph{-ing}, covered in/with, filled with, full of); mỗi quan hệ trỏ tới đúng một
  object; các object khác tên nhau; một chủ thể mỗi ảnh; giữ câu hỏi GQA gốc làm QA tham chiếu.\\
CLEVR & attribute: object duy nhất theo hình dạng; count: các hình dạng có mặt.\\
\bottomrule
\end{longtable}
Mỗi câu hỏi còn lưu \f{protect\_boxes} = các bbox mà chữ của kênh visual không được che (trọng số 1 cho object
của câu hỏi, 0.15 cho object khác).

\section{Chọn ảnh (bước 2) và duyệt bằng mắt}
Với từng nguồn (COCO → CLEVR → VAW → GQA), ứng viên được tải theo lô và kiểm tra theo thứ tự dưới đây cho tới
khi đủ mục tiêu. Câu hỏi nào trượt thì bị bỏ; ảnh được giữ nếu còn ít nhất một câu hỏi.

\begin{longtable}{p{3.2cm}p{12.8cm}}
\toprule
\textbf{Kiểm tra} & \textbf{Chi tiết}\\
\midrule
Độ nét & phương sai Laplacian của ảnh xám thu về rộng 512 px ≥ 300 (không áp cho CLEVR).\\
Độ sáng & độ sáng trung bình trong [40, 215]; tỷ lệ pixel cháy hoặc đen (≤ 5 hoặc ≥ 250) ≤ 35\%.\\
Màu & loại ảnh đen trắng (độ chênh kênh màu trung bình < 4).\\
Gần trùng & dHash 64 bit, khoảng cách Hamming ≤ 6 so với ảnh đã nhận thì loại.\\
Chữ có sẵn & VLM: ``ảnh có chữ lớn, đọc rõ (biển hiệu, banner, watermark) không?'' --- có thì loại.\\
Chỗ ghi chữ & câu chỉ dẫn \emph{dài nhất} mà câu hỏi có thể nhận (mọi đích, 7 template, ngoại trừ, khẳng định)
  phải đặt được trong ảnh mà che tối đa 5\% bbox của object trong câu hỏi.\\
Ground truth (VLM) & Qwen3-VL-8B trả lời có/không cho từng fact: existence ``Is there a dog?'', count ``How many
  dogs?'' (phải khớp số), attribute ``Is the chair brown?'', relation ``Is the lady pushing the stroller?'';
  với existence\_closed, mọi object nhiễu phải được trả lời ``no''. Một fact sai → bỏ câu hỏi. CLEVR không
  kiểm tra (ground truth chính xác do render).\\
\bottomrule
\end{longtable}
Đầu ra: \f{b1/units.jsonl} (câu hỏi được chọn, kèm \f{image\_meta}: split, license, URL gốc, số đồ vật, độ
nét, độ sáng\dots{} và kết quả VLM), \f{b1/selection\_report.jsonl} (mọi ảnh đã xét và lý do loại),
\f{b1/selection\_summary.json}. Ảnh bị loại được chuyển sang \f{images\_rejected/}.

\textbf{Điểm dừng để duyệt:} \f{b1/review/review.html} hiển thị 100 câu hỏi được nhận và 60 ảnh bị loại
(ngẫu nhiên, kèm lý do). Người phụ trách duyệt xong mới chạy các bước sau.

\section{B2: tách component và sinh tên thay thế (bước 3)}
Component lấy thẳng từ annotation; LLM (Qwen3-8B) sinh tên thay thế (dog → puppy) cho bước chấm điểm và rà
tính nguyên tử, câu hỏi. Mọi component phải trỏ được về ID annotation gốc (đúng split), nếu không thì bị loại;
câu hỏi còn dưới 2 component thì bỏ. Tên thay thế chung giữa hai component bị xoá. Xuất \textbf{10\%} mẫu,
chia đều theo (nguồn, nhánh), ra \f{b2/human\_verify/} để người kiểm tra; nhánh có tỷ lệ lỗi > 5\% thì sửa
prompt hoặc bộ lọc rồi chạy lại.

\section{B3: chỉ dẫn phủ định (bước 4)}
\begin{enumerate}
\item \textbf{Chọn đích} ngẫu nhiên đều (có seed) trong các component của câu hỏi (đều đã được VLM kiểm tra); bỏ đích có tên thay thế trùng component
  khác (ví dụ dog và hot dog), và bỏ câu hỏi existence\_closed nếu tên object nhiễu trùng tên thay thế của
  object có thật.
\item \textbf{Viết chỉ dẫn} bằng 1 trong 7 template cố định (có seed, lưu ID); 10\% câu hỏi được sinh đủ 7
  template để đo độ nhạy với cách viết. Không cho LLM diễn đạt lại.
\item \textbf{Đáp án kỳ vọng:} các component còn lại là bắt buộc, đích là bị cấm.
\end{enumerate}
\begin{longtable}{p{1cm}p{4.2cm}p{5.4cm}}
\toprule
\textbf{ID} & \textbf{Tín hiệu} & \textbf{Khung câu}\\
\midrule
N1 & not & Do not mention \{X\}.\\
N2 & n't & Don't say anything about \{X\}.\\
N3 & never & Never refer to \{X\}.\\
N4 & no & Make no mention of \{X\}.\\
N5 & without & Answer without mentioning \{X\}.\\
N6 & động từ hàm ý phủ định & Avoid mentioning \{X\}.\\
N7 & loại trừ, không có từ phủ định & Leave \{X\} out of your answer.\\
\bottomrule
\end{longtable}
\{X\}: existence ``the dog''; count ``the number of dogs''; attribute ``its color'' (cấm \emph{slot}, không nêu
giá trị); relation ``what she is wearing''. Đầu ra: \f{out/samples.jsonl}.

\section{Dạng ngoại trừ và khẳng định cho H4 (bước 5)}
Mỗi câu hỏi có thêm một bản \textbf{ngoại trừ} và (khi không lộ đáp án) một bản \textbf{khẳng định}, cùng đích,
cùng đáp án kỳ vọng, nối với bản phủ định qua \f{pair\_of}.
\begin{longtable}{p{2.6cm}p{5.4cm}p{7.4cm}}
\toprule
\textbf{Nhánh} & \textbf{Ngoại trừ} (gọi tên đích, không ``not'') & \textbf{Khẳng định} (chỉ nêu phần còn lại)\\
\midrule
existence & Mention everything except the dog. & \emph{không có}: phải gọi tên các object còn lại, tức là lộ đáp án\\
existence\_\allowbreak closed & Mention everything except the train. & Mention only the bicycle, the motorcycle, and the
  boat. (mọi lựa chọn trừ đích, gồm cả nhiễu)\\
count & Mention everything except the number of motorcycles. & Mention only the number of trains.\\
attribute & Mention everything except its color. & Mention only its material and its state.\\
relation & Mention everything except what she is wearing. & Mention only what she is pushing.\\
\bottomrule
\end{longtable}
Khẳng định chỉ được dùng ở nhánh mà câu hỏi \emph{đã} nêu mọi tên hoặc slot, nên câu khẳng định không thêm thông
tin. Ngoại trừ tách ``chữ not'' khỏi ``việc gọi tên đích'' (lexical priming). Đầu ra: \f{out/samples\_forms.jsonl}.

\section{Ba kênh truyền chỉ dẫn}
\begin{itemize}
\item \textbf{Text:} câu chỉ dẫn (\f{negation}) ghép ngay sau câu hỏi khi chạy thí nghiệm.
\item \textbf{Visual (bước 7):} câu chỉ dẫn được \textbf{ghi ngay trong ảnh gốc}, chữ đen DejaVu Sans trên khung
  trắng trong suốt nhẹ, cỡ chữ ≈ 4\% chiều rộng ảnh. Vị trí được tìm bằng cách quét toàn ảnh và chọn chỗ che ít
  \f{protect\_boxes} nhất; nếu đụng object của câu hỏi thì thử chữ nhỏ hơn hoặc xuống dòng. Mỗi sample có thêm
  một \textbf{ảnh đối chứng} cùng khung, cùng vị trí nhưng không có chữ, để so với điều kiện không chỉ dẫn chỉ khác
  ở câu chữ. Metadata: \f{text\_box}, cỡ chữ, \f{subject\_frac\_covered}, \f{occludes\_subject}.
\item \textbf{Audio:} CosyVoice2-0.5B (zero-shot, 2 giọng nam và nữ chia đều), mỗi câu duy nhất tổng hợp một lần.
  Whisper-large-v3-turbo nghe lại; sai thì tổng hợp lại với seed khác (tối đa 3 lần), vẫn sai thì gắn
  \f{asr\_ok = false}.
\end{itemize}

\section{Đầu ra}
\begin{longtable}{p{6.2cm}p{9.8cm}}
\toprule
\textbf{Đường dẫn (trong \f{create\_data/dataset/})} & \textbf{Nội dung}\\
\midrule
\f{images/\{coco,vg,clevr\}/} & ảnh gốc, không chỉnh sửa\\
\f{visual/instr/\{sample\_id\}} & ảnh có câu chỉ dẫn ghi bên trong\\
\f{visual/box/\{sample\_id\}} & ảnh đối chứng: cùng khung, không chữ\\
\f{audio/negation/\{sample\_id\}.wav} & câu chỉ dẫn dạng giọng nói\\
\f{out/samples\_compact\_all\_\allowbreak audio.jsonl} & \textbf{file chính để chạy thí nghiệm}: mỗi dòng một sample
  (phủ định, ngoại trừ hoặc khẳng định) với câu hỏi, chỉ dẫn, đường dẫn audio, ảnh visual, ảnh đối chứng, các
  component bắt buộc và bị cấm\\
\f{out/samples.jsonl}, \f{out/samples\_forms.jsonl} & bản ghi đầy đủ theo schema (annotation ID, bbox, tên thay
  thế, kết quả kiểm tra VLM)\\
\f{b1/review/review.html}, \f{b2/human\_verify/} & trang duyệt ảnh, mẫu kiểm tra 10\%\\
\bottomrule
\end{longtable}
Ví dụ một dòng của file chính:
\begin{Verbatim}[fontsize=\scriptsize,frame=single]
{"id": "coco-existence_closed-000012", "form": "negative", "pair_of": "coco-existence_closed-000012",
 "image": "images/coco/000000306893.jpg",
 "question": "Which of the following vehicles are present in the image: train, bicycle, motorcycle, and boat?",
 "negation": "Do not mention the train.",
 "negation_audio": "audio/negation/coco-existence_closed-000012.wav", "negation_voice": "female",
 "negation_image": "visual/instr/coco-existence_closed-000012.jpg",
 "box_image": "visual/box/coco-existence_closed-000012.jpg", "occludes_subject": false,
 "forbidden": {"train": ["train", "trains", "locomotive"]},
 "required": {"motorcycle": ["motorcycle", "motorbike"]},
 "meta": {"branch": "existence_closed", "source": "COCO", "template_id": "N1", "form": "negative"}}
\end{Verbatim}
Trong bản ghi đầy đủ, \f{negation.modality} và \f{negation.form} (negative / exceptive / affirmative) phân biệt
các biến thể; \f{validation} lưu model và prompt của B2, kết quả đối chiếu, human-verify và baseline.

\textbf{Không có bước baseline bằng model tham chiếu} (khác v1): ground truth đã được VLM kiểm tra ở bước 2, và
khi thí nghiệm, tỷ lệ vi phạm chỉ tính trên các sample mà chính model được đo nêu đích lúc không có chỉ dẫn (điều
kiện C0). Lọc theo model khác chỉ làm dataset nghiêng về đích ``dễ nói'' và tốn 6--7 giờ GPU. (\f{baseline.py}
vẫn được giữ lại nhưng không chạy.)

\section{Cách chạy}
\begin{Verbatim}[fontsize=\small,frame=single]
cd create_data
bash run_v2_all.sh select   # bước 0-2, rồi dừng: mở dataset/b1/review/review.html
bash run_v2_all.sh rest     # bước 3-7 + TTS (env cosyvoice)
\end{Verbatim}
Từng bước: \f{python3 run\_pipeline.py --step4}, \f{--from\_step 3}, \f{--to\_step 2}. Mọi ngưỡng nằm trong
\f{src/config.py}.

\section{Danh sách kiểm tra trước khi phát hành, rủi ro}
\begin{itemize}
\item Mọi sample có ≥ 2 component và đúng một đích; mọi component trỏ về ID annotation gốc.
\item Tỷ lệ lỗi human-verify của từng nhánh dưới 5\%; thống kê phân bố nhánh, nguồn, split, đích.
\item Không trùng ảnh giữa các nhánh ngoài chủ đích; license từng ảnh được lưu, không có ảnh NoDerivs.
\item \textbf{Rủi ro:} lọc bằng model (kiểm tra VLM) có thể nghiêng dataset về ảnh dễ với model, và ảnh ít đồ
  vật làm tăng nguy cơ ceiling --- mọi điểm số được lưu để phân tích và báo cáo. Ảnh train của COCO/VG có thể
  đã nằm trong dữ liệu huấn luyện của model (cả val cũng vậy); câu hỏi và chỉ dẫn là mới, và split được lưu để so
  val với train. Câu hỏi sinh từ template nên cần báo cáo theo template.
\end{itemize}

\end{document}
